% =====================================================================
%  Bien dich: pdfLaTeX (mac dinh cua Overleaf) -- khong can chinh gi.
%  Neu muon dung font he thong thi doi compiler sang XeLaTeX va thay
%  ba dong inputenc/fontenc/lmodern ben duoi bang:
%      \usepackage{fontspec}
%      \setmainfont{TeX Gyre Termes}
% =====================================================================
\documentclass[11pt,a4paper]{article}

\usepackage[utf8]{inputenc}
\usepackage[T5]{fontenc}
\usepackage{lmodern}

\usepackage[margin=2.5cm]{geometry}
\usepackage{amsmath}
\usepackage{amssymb}
\usepackage{booktabs}
\usepackage{array}
\usepackage{caption}
\usepackage{microtype}
\usepackage{enumitem}
\usepackage{framed}
\usepackage{tikz}
\usetikzlibrary{arrows.meta, positioning}
\usepackage[hidelinks]{hyperref}

\renewcommand{\tablename}{Bảng}
\renewcommand{\figurename}{Hình}
\renewcommand{\contentsname}{Nội dung}
\renewcommand{\refname}{Tài liệu tham khảo}

\captionsetup[table]{skip=6pt}
\captionsetup[figure]{skip=10pt}
\setlist[itemize]{itemsep=3pt, topsep=4pt}
\setlist[enumerate]{itemsep=4pt, topsep=4pt}

\newcolumntype{L}{>{\raggedright\arraybackslash}p{6.6cm}}
\newcolumntype{M}{>{\raggedright\arraybackslash}p{8.6cm}}
\newenvironment{luuy}{\begin{leftbar}\noindent\ignorespaces}{\end{leftbar}}

\title{\bfseries Đọc hiểu RanPAC\\[2pt]
\large\mdseries Chiếu ngẫu nhiên và mô hình tiền huấn luyện cho học liên tục}
\author{Tóm tắt và diễn giải bài báo của McDonnell và cộng sự, NeurIPS 2023}
\date{Tháng 8, 2026}

\begin{document}
\maketitle

\noindent
Tài liệu này diễn giải bài báo RanPAC \cite{ranpac} cho người mới bắt đầu: bài
giải quyết vấn đề gì, ý tưởng cốt lõi là gì, từng thành phần hoạt động ra sao, và
vì sao nó liên quan tới đề tài của chúng ta. Mọi số liệu đều lấy trực tiếp từ bài
báo, có ghi rõ nguồn là bảng nào.

\tableofcontents
\bigskip

% =====================================================================
\section{Bài báo giải quyết vấn đề gì}

\subsection{Học liên tục và quên thảm khốc}

Bối cảnh là \textbf{học liên tục} (continual learning): dữ liệu đến theo từng
đợt, mô hình phải học đợt mới mà không được xem lại đợt cũ, và không được lưu ảnh
cũ. Trở ngại trung tâm là \textbf{quên thảm khốc} -- việc cập nhật trọng số cho
dữ liệu mới sẽ ghi đè lên những gì đã học trước đó.

\subsection{Hai giao thức: CIL và DIL}

Bài báo làm việc với hai kịch bản khác nhau, và cần phân biệt rõ vì kết quả được
báo cáo riêng cho từng loại.

\begin{itemize}[leftmargin=*]
\item \textbf{Class-Incremental Learning (CIL).} Mỗi đợt dạy một tập lớp
      \emph{khác nhau}, không giao nhau. Ví dụ đợt 1 học lớp 0--9, đợt 2 học lớp
      10--19. Điểm khó: lúc kiểm tra, mô hình \emph{không được biết} ảnh thuộc
      đợt nào, phải chọn đúng lớp trong toàn bộ các lớp đã gặp.
\item \textbf{Domain-Incremental Learning (DIL).} Mọi đợt dùng \emph{cùng một tập
      lớp}, nhưng dữ liệu mỗi đợt đến từ một miền khác nhau. Ví dụ đợt 1 là ảnh
      chụp, đợt 2 là tranh vẽ, đợt 3 là ảnh phác thảo.
\end{itemize}

Đề tài của chúng ta thuộc loại thứ nhất, CIL.

\subsection{Ba hướng khai thác mô hình tiền huấn luyện}

Từ khi có các mô hình nền tảng lớn được tiền huấn luyện sẵn, lĩnh vực này tách
thành ba hướng.

\begin{table}[htbp]
\centering
\caption{Ba hướng khai thác mô hình tiền huấn luyện cho học liên tục.}
\label{tab:huong}
\begin{tabular}{lLL}
\toprule
Hướng & Cách làm & Điểm yếu \\
\midrule
Prompting
  & Học một túi prompt nhỏ gắn vào đầu vào của transformer (L2P, DualPrompt,
    CODA-Prompt)
  & Chỉ áp dụng được cho transformer \\
\addlinespace
Tinh chỉnh backbone
  & Tinh chỉnh nhẹ toàn mạng với tốc độ học thấp (SLCA)
  & Vẫn quên tích lũy, tốn tài nguyên huấn luyện \\
\addlinespace
Class-Prototype (CP)
  & Lấy trung bình đặc trưng của mỗi lớp làm ``đại diện'', phân loại bằng độ
    tương đồng cosine (ADaM)
  & Bị coi là đơn giản nhưng kém chính xác \\
\bottomrule
\end{tabular}
\end{table}

RanPAC thuộc hướng thứ ba. Luận điểm của bài là: \textbf{hướng CP bị đánh giá
thấp oan -- không phải vì bản chất nó yếu, mà vì người ta đang dùng nó sai cách.}

\subsection{Ký hiệu dùng trong tài liệu}

Bảng này gom mọi ký hiệu sẽ gặp về một chỗ, để không phải lần ngược tìm nghĩa.

\begin{table}[htbp]
\centering
\caption{Ký hiệu. Các giá trị số là cấu hình chính của bài báo.}
\label{tab:kyhieu}
\begin{tabular}{lM}
\toprule
Ký hiệu & Ý nghĩa \\
\midrule
$T$        & Số tác vụ trong chuỗi (thường là 10) \\
$D_t$      & Dữ liệu huấn luyện của tác vụ thứ $t$ \\
$N_t$, $N$ & Số mẫu trong tác vụ $t$, và tổng số mẫu \\
$K$        & Tổng số lớp sau khi học hết $T$ tác vụ \\
$L$        & Số chiều đặc trưng gốc do backbone sinh ra ($L = 768$) \\
$M$        & Số chiều sau khi chiếu ngẫu nhiên ($M = 10\,000$) \\
$f$        & Vector đặc trưng do backbone đóng băng sinh ra, $L$ chiều \\
$W$        & Ma trận chiếu ngẫu nhiên $L \times M$, đóng băng, không huấn luyện \\
$\phi$     & Hàm phi tuyến áp theo từng phần tử (bài dùng ReLU) \\
$h$        & Đặc trưng sau chiếu và kích hoạt, $h = \phi(f^{\top}W)$, $M$ chiều \\
$H$        & Ma trận gom mọi $h$ của toàn bộ dữ liệu đã thấy \\
$c_y$      & Class-prototype của lớp $y$, tính trong không gian $M$ chiều \\
$C$        & Ma trận gom mọi $c_y$ thành các cột \\
$G$        & Ma trận Gram, $G = HH^{\top}$, kích thước $M \times M$ \\
$\lambda$  & Hệ số điều chuẩn của hồi quy ridge \\
$W_o$      & Trọng số phân lớp cuối cùng, $W_o = (G + \lambda I)^{-1}C$ \\
\bottomrule
\end{tabular}
\end{table}

% =====================================================================
\section{Chẩn đoán: vì sao class-prototype đang kém}

Đây là phần thuyết phục nhất của bài, vì họ không suy đoán mà đo trực tiếp.

Lấy ViT-B/16 tiền huấn luyện trên ImageNet-21K, chạy trên ImageNet-R, rồi tính
\textbf{hệ số tương quan Pearson giữa từng cặp class-prototype}. Kết quả: các
prototype của những lớp \emph{khác nhau} lại tương quan với nhau rất mạnh, cao
hơn hai lần so với khi làm y hệt trên chính ImageNet-1K.

Diễn giải: khi đem mô hình sang một miền dữ liệu mới, \textbf{mọi đặc trưng đều
bị dồn về cùng một hướng}. Prototype của hai lớp hoàn toàn khác nhau trở nên
giống nhau một cách bất thường.

Hậu quả đo được ngay trên histogram độ tương đồng: phân bố ``giống lớp đúng của
mình'' và ``giống các lớp sai'' chồng lấn nặng.

\begin{table}[htbp]
\centering
\caption{Độ chính xác trên tập huấn luyện ImageNet-R, theo Hình 2 của bài báo.
Con số cho thấy vấn đề nằm ở cách chấm điểm, không ở bản thân ý tưởng prototype.}
\label{tab:chandoan}
\begin{tabular}{lr}
\toprule
Cách chấm điểm & Độ chính xác huấn luyện \\
\midrule
Class-prototype kèm cosine (NCM)                     & 64\,\% \\
Class-prototype kèm nghịch đảo Gram (Mục \ref{sec:gram})
                                                     & 75\,\% \\
Linear probe huấn luyện chung trên cả 200 lớp        & 82\,\% \\
Gram kèm chiếu ngẫu nhiên, $M = 2000$ (Mục \ref{sec:rp})
                                                     & \textbf{87\,\%} \\
\bottomrule
\end{tabular}
\end{table}

Hàng cuối đáng chú ý: nó \emph{vượt} cả linear probe được huấn luyện chung trên
toàn bộ dữ liệu. Hai mục tiếp theo giải thích hai bước đưa 64\,\% lên 87\,\%.

% =====================================================================
\section{Ý tưởng thứ nhất: khử tương quan bằng ma trận Gram}
\label{sec:gram}

\subsection{Đổi công thức chấm điểm}

Cách cũ (NCM) dùng cosine giữa đặc trưng và prototype \emph{trung bình}
$\bar c_y$:
\begin{equation}
s_y = \frac{f_{\text{test}}^{\top} \bar c_y}
           {\|f_{\text{test}}\| \cdot \|\bar c_y\|}.
\label{eq:cosine}
\end{equation}
RanPAC thay bằng nghịch đảo ma trận Gram:
\begin{equation}
s_y = f_{\text{test}}^{\top} G^{-1} c_y .
\label{eq:gram}
\end{equation}

Chú ý một chi tiết nhỏ dễ bỏ qua: ở đây $c_y$ là \emph{tổng} các đặc trưng của
lớp $y$ chứ không phải trung bình. Bỏ phép chia cho số mẫu không gây hại, vì
nghịch đảo $G^{-1}$ và hệ số $\lambda$ đã tự điều chỉnh thang đo; giữ dạng tổng
thì việc cộng dồn qua các tác vụ đơn giản hơn nhiều.

\subsection{Trực giác: tẩy trắng rồi mới so sánh}

Ta đi tìm một phép biến đổi tuyến tính $D$ sao cho sau khi biến đổi, ma trận Gram
trở thành ma trận đơn vị -- nghĩa là các chiều đặc trưng hết tương quan với nhau:
\begin{equation}
(DH)(DH)^{\top} = I \;\Longrightarrow\; D^{\top}D = (HH^{\top})^{-1} = G^{-1}.
\end{equation}
Khi đó tích vô hướng giữa đặc trưng đã tẩy trắng và prototype đã tẩy trắng chính
là công thức \eqref{eq:gram}:
\begin{equation}
\hat s_y = (D h_{\text{test}})^{\top}(D c_y)
         = h_{\text{test}}^{\top} G^{-1} c_y .
\end{equation}

\begin{luuy}
Nói ngắn gọn: \textbf{tẩy trắng dữ liệu rồi mới so sánh}, thay vì so sánh trực
tiếp trên một không gian đã bị méo. Đó là toàn bộ nội dung của $G^{-1}$.
\end{luuy}

\subsection{Ba mối liên hệ lý thuyết đáng nhớ}

\paragraph{Liên hệ với hồi quy ridge.}
Dạng $W_o = (G + \lambda I)^{-1} C$ \textbf{chính xác là nghiệm đóng} của bài
toán bình phương tối thiểu có điều chuẩn $l_2$, khớp nhãn one-hot:
\begin{equation}
W_o = \arg\min_{W}
\left( \|Y_{\text{train}}^{\top} - W^{\top}H\|_2^2 + \lambda\|W\|_2^2 \right).
\end{equation}
Hệ quả thực tiễn rất mạnh: huấn luyện một lớp tuyến tính bằng SGD với mất mát
bình phương và weight decay bằng $\lambda$ thì \emph{giỏi lắm cũng chỉ tiệm cận}
được nghiệm này. Ở đây ta tính thẳng bằng công thức, không cần lặp.

\paragraph{Liên hệ với LDA và khoảng cách Mahalanobis.}
Công thức \eqref{eq:gram} rất giống LDA, nhưng LDA dùng ma trận \emph{hiệp phương
sai} $S$ còn RanPAC dùng ma trận \emph{Gram} $G$. Hai cái chỉ trùng nhau khi mọi
chiều đặc trưng có trung bình bằng $0$, điều nói chung không đúng. Dùng $G$ có
hai cái lợi: không phải tính thành phần bias, và việc cộng dồn qua các tác vụ đơn
giản hơn vì không phải trừ trung bình.

\paragraph{Liên hệ với ZCA whitening.}
Nếu tất cả các lớp có xác suất tiên nghiệm bằng nhau, phân loại LDA tương đương
với việc tẩy trắng ZCA rồi tìm khoảng cách Euclid nhỏ nhất. Công thức
\eqref{eq:gram} cùng tinh thần đó, chỉ khác ở chỗ đưa ma trận Gram về đơn vị thay
vì đưa hiệp phương sai về đơn vị.

% =====================================================================
\section{Ý tưởng thứ hai: chiếu ngẫu nhiên kèm phi tuyến}
\label{sec:rp}

Chiếu đặc trưng $L = 768$ chiều lên $M = 10\,000$ chiều bằng ma trận ngẫu nhiên
$W$ đóng băng, rồi cho qua hàm phi tuyến $\phi$ (ReLU):
\begin{equation}
h = \phi(f^{\top} W).
\end{equation}

\subsection{Vì sao phi tuyến là bắt buộc}

Đây là điểm dễ hiểu nhầm nhất, và lập luận rất gọn:

\begin{luuy}
Nếu bỏ $\phi$ đi, phép chiếu ngẫu nhiên chỉ còn là một \textbf{phép đổi cơ sở
tuyến tính}. Mà một bộ phân lớp tuyến tính thì \emph{bất biến} với phép đổi cơ
sở. Chiếu lên một triệu chiều cũng vô ích.
\end{luuy}

Thí nghiệm ở Hình 3(a) của bài xác nhận đúng như vậy: bỏ kích hoạt phi tuyến thì
dù $M$ lớn đến đâu, kết quả cũng không hơn gì không chiếu.

\subsection{Vậy phi tuyến mang lại gì}

Nó tạo ra \textbf{số hạng tương tác giữa các đặc trưng}. Khai triển Taylor của
$\phi$ sinh ra các tích chéo giữa những chiều đã chiếu -- thứ mà một mô hình
tuyến tính trên đặc trưng gốc không bao giờ chạm tới được.

Bài chứng minh điều này bằng một thí nghiệm rất sạch (Hình 3(b)): lấy 100 đặc
trưng đầu tiên, dựng \emph{tường minh} toàn bộ tích chéo từng cặp, rồi phân loại.
Kết quả tốt hơn hẳn đặc trưng thô. Nhưng với $L = 768$ thì số tích chéo tường
minh lên tới gần $300\,000$, bất khả thi về tính toán. Chiếu ngẫu nhiên kèm phi
tuyến là \textbf{cách rẻ để lấy gần hết lợi ích đó}.

\subsection{Vì sao phải mở rộng chiều chứ không nén}

Có hai lý do, một thực nghiệm và một hình học.

\emph{Thực nghiệm:} vì $W$ là ngẫu nhiên chứ không được học, nên cần dư thừa để
bù lại. Bài thử nén từ 768 xuống 500 chiều và kết quả \emph{thấp hơn cả} phương
án không dùng chiếu.

\emph{Hình học:} Phụ lục B.1 và B.2 dùng bất đẳng thức Chernoff để chỉ ra rằng
khi $M$ tăng, chuẩn của các vector đã chiếu càng tập trung quanh kỳ vọng, tức là
phân bố càng tiến gần một Gaussian đẳng hướng. Điều đó quan trọng vì công thức
\eqref{eq:gram} ngầm giả định dữ liệu có dạng Gaussian. Đồng thời, $M$ càng lớn
thì tích vô hướng giữa hai mẫu bất kỳ càng dễ khác biệt nhau -- nghĩa là các lớp
càng dễ tách.

\begin{table}[htbp]
\centering
\caption{Độ chính xác trên Split-CIFAR100 theo số chiều chiếu $M$, không dùng
Phase 1, $\lambda$ cố định bằng 100 (Bảng A5 của bài báo).}
\label{tab:scale}
\begin{tabular}{lrrrrrrr}
\toprule
$M$ & 100 & 400 & 1250 & 2500 & 5000 & 10000 & 15000 \\
\midrule
Độ chính xác (\%) & 71.6 & 83.9 & 86.8 & 87.7 & 88.4 & 88.8 & 89.0 \\
\bottomrule
\end{tabular}
\end{table}

Mức tăng bão hòa dần sau khoảng $M = 5000$. Bài cũng lưu ý rằng $M$ phải lớn hơn
khoảng $1250$ thì mới vượt được phương án không dùng chiếu.

\subsection{Bản thân ý tưởng chiếu ngẫu nhiên không mới}

Cần nói rõ điều này để đánh giá đúng đóng góp. Chiếu ngẫu nhiên đóng băng kèm phi
tuyến đã được đề xuất và dùng nhiều lần từ những năm 1990 trong dòng nghiên cứu
về mạng nơ-ron ngẫu nhiên, nhưng luôn ở vai trò \emph{bộ học độc lập, không liên
tục}.

Đóng góp thật của RanPAC nằm ở hai quan sát nối chúng lại với học liên tục:

\begin{enumerate}[leftmargin=*]
\item Các chiến lược class-prototype hiện có cho học liên tục \emph{tương đương}
      với các mô hình tuyến tính học sau phép chiếu ngẫu nhiên trong dòng nghiên
      cứu cũ.
\item Trọng số đóng băng và không được huấn luyện thì \emph{không thể gây quên}.
\end{enumerate}

Nói cách khác, RanPAC không phát minh cơ chế mới; nó nhận ra một cơ chế cũ có
đúng tính chất mà học liên tục cần.

% =====================================================================
\section{Vì sao cách này hợp với học liên tục đến mức hoàn hảo}
\label{sec:coloi}

Đây là điểm cốt lõi của cả bài, và nó rất thanh lịch.

\begin{equation}
G = \sum_{t=1}^{T}\sum_{n=1}^{N_t} h_{t,n} \otimes h_{t,n},
\qquad
C = \sum_{t=1}^{T}\sum_{n=1}^{N_t} h_{t,n} \otimes y_{t,n}.
\label{eq:congdon}
\end{equation}

Cả hai đều là \textbf{tổng cộng dồn}. Mà phép cộng có tính giao hoán, nên:

\begin{luuy}
Học tuần tự qua $T$ tác vụ cho ra kết quả \textbf{giống hệt} học tất cả cùng một
lúc. Không phải ``quên ít'', mà là \textbf{không quên, chứng minh được bằng đại
số}, không cần thí nghiệm. Thứ tự các tác vụ hoàn toàn không ảnh hưởng.
\end{luuy}

Lý do sâu xa: không có bước cập nhật trọng số lặp nào để mà ghi đè kiến thức cũ.
Toàn bộ ``huấn luyện'' của Phase 2 chỉ là cộng ma trận.

Đó cũng là lý do bài nhấn mạnh cụm \emph{training-free}: khả năng chống quên
không đến từ một cơ chế chống quên nào cả, mà đến từ việc \textbf{chọn một dạng
bài toán vốn không có chỗ cho việc quên xảy ra}.

Một hệ quả kiểm chứng được, và bài có kiểm chứng ở Phụ lục F: khi bỏ Phase 1, kết
quả cuối cùng gần như \emph{không dao động} dù đổi cách gán ngẫu nhiên lớp vào
tác vụ, và cũng không đổi khi chia dữ liệu thành 5, 10 hay 20 tác vụ. Khi có
Phase 1, dao động tăng lên rõ rệt -- vì Phase 1 dùng SGD, mà SGD thì có ngẫu
nhiên.

% =====================================================================
\section{Phase 1: vá khoảng cách miền}

\subsection{Lý do tồn tại}

Vẫn còn một vấn đề. Đặc trưng đóng băng có thể không hợp với miền dữ liệu mới --
ảnh phong cách hóa, ảnh y tế, ảnh vệ tinh.

Phase 1 xử lý bằng cách huấn luyện một bộ \textbf{PETL} (Parameter-Efficient
Transfer Learning) \textbf{chỉ trên tác vụ đầu tiên} $D_1$, rồi đóng băng vĩnh
viễn. Đầu phân lớp tạm dùng để huấn luyện nó bị vứt bỏ trước khi vào Phase 2.

Logic gồm hai vế. Thứ nhất, chỉ huấn luyện ở một tác vụ duy nhất thì không có cơ
hội để quên tích lũy. Thứ hai, dữ liệu tác vụ 1 dù ít vẫn đại diện cho miền mới
tốt hơn nhiều so với ImageNet-21K.

\subsection{Ba phương pháp PETL được thử}

Cả ba đều \textbf{không đụng vào một trọng số nào của mạng gốc}, chỉ thêm tham số
mới rất nhỏ.

\begin{itemize}[leftmargin=*]
\item \textbf{VPT} (Visual Prompt Tuning): thêm một số token ``prompt'' học được
      vào đầu vào của các tầng transformer. Bài dùng bản \emph{deep}, tức là chèn
      prompt ở mọi tầng, với độ dài prompt bằng 5.
\item \textbf{AdaptFormer}: chèn một khối MLP cổ chai nhỏ song song với khối MLP
      sẵn có của transformer -- chiếu xuống chiều thấp, qua phi tuyến, rồi chiếu
      ngược lên. Bài dùng chiều chiếu bằng 64.
\item \textbf{SSF} (Scaling and Shifting Features): học một hệ số nhân và một hệ
      số cộng cho từng kênh đặc trưng. Đây là phương pháp ít tham số nhất trong
      ba.
\end{itemize}

Phụ lục F.6 cho thấy phương pháp nào tốt nhất \emph{thay đổi theo bộ dữ liệu và
theo backbone}, không có lựa chọn thắng tuyệt đối.

Chi tiết huấn luyện chung: SGD, batch 48, tốc độ học 0.01, weight decay 0.0005,
momentum 0.9, 20 epoch, lịch cosine annealing.

% =====================================================================
\section{Thuật toán đầy đủ}

\begin{figure}[htbp]
\centering
\resizebox{\textwidth}{!}{%
\begin{tikzpicture}[
  font=\footnotesize,
  >={Stealth[length=2.2mm]},
  box/.style={draw, align=center, text width=3.0cm, minimum height=15mm,
              inner sep=4pt},
  frz/.style={box, fill=black!8},
  ttl/.style={font=\footnotesize\bfseries, anchor=south west}
]
% ---------- Phase 1 ----------
\node[box] (d1) at (0,0) {Dữ liệu\\tác vụ đầu tiên};
\node[box, right=12mm of d1]   (petl) {Backbone\\ + tham số PETL};
\node[box, right=12mm of petl] (tmp)  {Đầu phân lớp tạm\\ (vứt bỏ sau đó)};
\draw[->] (d1)   -- (petl);
\draw[->] (petl) -- (tmp);
\node[ttl, yshift=3mm] at (d1.north west) {Phase 1 --- thích nghi phiên đầu};

% ---------- Phase 2 ----------
\node[box, below=34mm of d1]  (dt)  {Dữ liệu\\mọi tác vụ};
\node[frz, right=12mm of dt]  (bb)  {Backbone\\đã thích nghi\\(đóng băng)};
\node[frz, right=12mm of bb]  (rp)  {Chiếu ngẫu nhiên\\rồi ReLU\\(đóng băng)};
\node[box, right=12mm of rp]  (acc) {Cộng dồn\\$G$ và $C$};
\node[box, right=12mm of acc] (sol) {Giải ridge\\$(G+\lambda I)^{-1}C$};
\draw[->] (dt)  -- (bb);
\draw[->] (bb)  -- (rp);
\draw[->] (rp)  -- (acc);
\draw[->] (acc) -- (sol);
\node[ttl, yshift=3mm] at (dt.north west) {Phase 2 --- học liên tục};

% ---------- noi hai pha ----------
\draw[->, dashed] (petl) -- (bb);
\end{tikzpicture}}
\caption{Hai pha của RanPAC. Các khối tô xám là thành phần đóng băng, không có
tham số nào được huấn luyện trong đó. Mũi tên nét đứt là tham số PETL học được ở
Phase 1 rồi đóng băng và dùng lại cho Phase 2. Đặc trưng ra khỏi backbone có 768
chiều; sau khi chiếu và kích hoạt thì có $10\,000$ chiều. Toàn bộ ``huấn luyện''
của Phase 2 chỉ là cộng dồn hai ma trận rồi giải một hệ tuyến tính.}
\label{fig:ranpac}
\end{figure}

\paragraph{Phase 1 -- thích nghi phiên đầu (tùy chọn).}
\begin{enumerate}[leftmargin=*]
\item Với mỗi ảnh trong $D_1$, trích đặc trưng và huấn luyện tham số PETL bằng
      SGD với mất mát cross-entropy trên các lớp của tác vụ đầu.
\item Đóng băng tham số PETL. Vứt bỏ đầu phân lớp tạm.
\end{enumerate}

\paragraph{Phase 2 -- học liên tục.}
\begin{enumerate}[leftmargin=*]
\item Sinh ma trận chiếu $W \sim \mathcal{N}(0,1)$ kích thước $L \times M$, đóng
      băng vĩnh viễn.
\item Với mỗi tác vụ $t = 1,\ldots,T$ và mỗi ảnh trong $D_t$: trích đặc trưng
      $f$, tính $h = \phi(f^{\top}W)$, rồi cộng dồn $G$ và $C$ theo
      \eqref{eq:congdon}.
\item Cuối mỗi tác vụ: chọn $\lambda$, rồi tính $W_o = (G + \lambda I)^{-1}C$.
\end{enumerate}

\subsection{Cách chọn \texorpdfstring{$\lambda$}{lambda}}

Bước này đáng chú ý vì nó tuân thủ đúng ràng buộc của học liên tục. Ở mỗi tác vụ:

\begin{enumerate}[leftmargin=*]
\item Chia dữ liệu \emph{của chính tác vụ đó} theo tỉ lệ 80:20.
\item Quét $\lambda$ qua 17 bậc độ lớn, từ $10^{-8}$ đến $10^{8}$.
\item Với mỗi $\lambda$, tính $W_o$ từ phần 80\,\% rồi đo sai số bình phương
      trung bình trên phần 20\,\%.
\item Chọn $\lambda$ cho sai số nhỏ nhất.
\end{enumerate}

Không đụng tới dữ liệu của các tác vụ cũ ở bất kỳ bước nào. Bài cũng thẳng thắn
ghi nhận rằng chọn $\lambda$ chỉ dựa trên tác vụ hiện tại có thể chưa tối ưu so
với khi được xem toàn bộ dữ liệu.

% =====================================================================
\section{Kết quả}

\subsection{Hai chỉ số đánh giá}

Gọi $R_{t,i}$ là độ chính xác trên tác vụ $i$ sau khi đã học xong tác vụ $t$.

\begin{equation}
A_t = \frac{1}{t}\sum_{i=1}^{t} R_{t,i},
\qquad
F_t = \frac{1}{t-1}\sum_{i=1}^{t-1}
      \max_{t' \in \{1,\ldots,t-1\}} \left( R_{t',i} - R_{t,i} \right).
\label{eq:chiso}
\end{equation}

$A_t$ là \textbf{Average Accuracy}, trung bình độ chính xác trên mọi tác vụ đã
học, đo tại giai đoạn $t$. Bài báo báo cáo $A_T$, tức giá trị ở giai đoạn cuối
cùng. $F_t$ là \textbf{Average Forgetting}, đo mức tụt của mỗi tác vụ so với đỉnh
hiệu năng của chính nó; giá trị càng thấp càng tốt.

\begin{luuy}
$A_T$ của bài báo \textbf{trùng định nghĩa} với cột Acc@1 trong bảng kết quả của
chúng ta. Đây là điều cần biết trước khi đối chiếu số liệu giữa hai bên.
\end{luuy}

\subsection{Các bộ dữ liệu}

\begin{table}[htbp]
\centering
\caption{Bảy bộ dữ liệu CIL (Bảng A2 của bài báo). $N$ là số mẫu huấn luyện, $K$
là tổng số lớp.}
\label{tab:dulieu}
\begin{tabular}{lrrrL}
\toprule
Bộ dữ liệu & $N$ & Số mẫu kiểm tra & $K$ & Nội dung \\
\midrule
CIFAR100      & 50\,000 & 10\,000 & 100 & Ảnh tự nhiên cỡ nhỏ, phóng lên 224 \\
ImageNet-R    & 24\,000 & 6\,000  & 200 & ``Rendition'' của 200 lớp ImageNet: tranh, hoạt hình, phác thảo, origami. Cố tình xa phân bố tiền huấn luyện \\
ImageNet-A    & 5\,981  & 5\,985  & 200 & Ảnh thật mà các bộ phân lớp ImageNet chuẩn hay đoán sai \\
CUB           & 9\,430  & 2\,358  & 200 & 200 loài chim, phân loại mịn \\
OmniBenchmark & 89\,697 & 5\,985  & 300 & Trải rộng nhiều lĩnh vực thị giác \\
VTAB          & 1\,796  & 8\,619  & 50  & 5 tác vụ đến từ các miền khác nhau; mang đặc điểm của cả CIL lẫn DIL \\
Stanford Cars & 8\,144  & 8\,041  & 196 & 196 mẫu xe hơi, phân loại mịn \\
\bottomrule
\end{tabular}
\end{table}

Ba bộ DIL là CORe50 (8 giai đoạn), CDDB-Hard (5 miền, bài toán nhị phân phát hiện
ảnh giả) và DomainNet (6 miền: ảnh thật, quickdraw, tranh vẽ, phác thảo, đồ họa
thông tin, clipart).

\begin{luuy}
Bảng này giải thích một điều quan trọng cho đề tài của chúng ta: \textbf{CUB là
ảnh chim, mà ImageNet-21K vốn rất giàu lớp chim}, nên đặc trưng đóng băng đã đủ
tốt ở đó. ImageNet-R thì ngược lại -- ảnh phong cách hóa, xa hẳn dữ liệu tiền
huấn luyện, nên bắt buộc phải có thích nghi theo miền. Đó chính là lý do các
phương pháp không huấn luyện thường chỉ thắng trên CUB.
\end{luuy}

\subsection{Kết quả class-incremental}

\begin{table}[htbp]
\centering
\caption{Kết quả và khảo sát loại bỏ thành phần cho CIL, $T = 10$ (Bảng 1 và
Bảng 2 của bài báo). Hai hàng in nghiêng là mốc tham chiếu, không phải phương
pháp học liên tục.}
\label{tab:ketqua}
\begin{tabular}{lrrrr}
\toprule
Cấu hình & CIFAR100 & ImageNet-R & ImageNet-A & CUB \\
\midrule
NCM đơn thuần (cosine)              & 83.4 & 61.2 & 49.3 & 85.1 \\
NCM kèm Phase 1                     & 87.8 & 70.1 & 49.7 & 85.4 \\
Gram, không chiếu, không Phase 1    & 87.0 & 67.6 & 54.4 & 86.8 \\
Gram kèm chiếu, không Phase 1       & 89.0 & 71.8 & 58.2 & 88.7 \\
Gram kèm Phase 1, không chiếu       & 90.6 & 73.9 & 57.7 & 86.6 \\
\midrule
\textbf{RanPAC đầy đủ}              & \textbf{92.2} & \textbf{77.9} & \textbf{62.4} & \textbf{90.3} \\
\midrule
\emph{Linear probe huấn luyện chung}       & \emph{87.9} & \emph{72.0} & \emph{56.6} & \emph{88.7} \\
\emph{Fine-tune toàn bộ, huấn luyện chung} & \emph{93.8} & \emph{86.6} & \emph{70.8} & \emph{90.5} \\
\bottomrule
\end{tabular}
\end{table}

Ba điều đọc ra được:

\begin{enumerate}[leftmargin=*]
\item \textbf{Chiếu ngẫu nhiên và Phase 1 bổ trợ nhau, không thay thế nhau.} Bỏ
      chiếu mất 1.6--4 điểm; bỏ Phase 1 mất 2.2--6.1 điểm; bỏ cả hai mất
      5.2--10.3 điểm.
\item \textbf{RanPAC vượt cả linear probe được huấn luyện chung} ở cả bảy bộ dữ
      liệu trong bài. Đây là điều đáng ngạc nhiên: một phương pháp học tuần tự
      lại hơn một mô hình được xem toàn bộ dữ liệu cùng lúc. Lý do là linear probe
      làm việc trên 768 chiều, còn RanPAC làm việc trên $10\,000$ chiều có tương
      tác phi tuyến.
\item \textbf{Khoảng cách tới cận trên rất hẹp trên một số tập.} Trên CIFAR100 nó
      chỉ kém fine-tune toàn bộ 1.6 điểm, trên CUB kém 0.2 điểm. Khoảng cách lớn
      chỉ còn ở hai biến thể ImageNet.
\end{enumerate}

\subsection{Kết quả domain-incremental}

\begin{table}[htbp]
\centering
\caption{Kết quả cho DIL (Bảng 3 của bài báo).}
\label{tab:dil}
\begin{tabular}{lrrr}
\toprule
Cấu hình & CORe50 & CDDB-Hard & DomainNet \\
\midrule
L2P                                & 78.3 & 61.3 & 40.2 \\
S-iPrompts                         & 83.1 & 74.5 & 50.6 \\
ADaM                               & 92.0 & 70.7 & 50.3 \\
\midrule
\textbf{RanPAC đầy đủ}             & \textbf{96.7} & \textbf{86.2} & \textbf{66.6} \\
\midrule
NCM đơn thuần                      & 80.4 & 69.8 & 46.4 \\
Gram, không chiếu, không Phase 1   & 92.4 & 80.7 & 54.1 \\
Gram kèm chiếu, không Phase 1      & 94.3 & 81.6 & 64.7 \\
Gram kèm Phase 1, không chiếu      & 94.8 & 84.2 & 55.2 \\
\bottomrule
\end{tabular}
\end{table}

Xu hướng giống hệt CIL, nhưng có một khác biệt đáng chú ý: trên DomainNet, phép
chiếu ngẫu nhiên đóng góp rất lớn còn Phase 1 gần như không thêm gì (64.7 lên
66.6). Điều này hợp lý với bản chất của DIL -- Phase 1 chỉ được huấn luyện trên
miền đầu tiên, nên nó không giúp được cho các miền sau vốn khác hẳn.

\subsection{Ảnh hưởng của số tác vụ}
\label{sec:sotacvu}

\begin{table}[htbp]
\centering
\caption{Kết quả theo số tác vụ $T$, dùng AdaptFormer cho Phase 1 (Bảng A4 của
bài báo). Cột cuối không phụ thuộc $T$, đúng như tính chất bất biến ở Mục
\ref{sec:coloi}.}
\label{tab:sotacvu}
\begin{tabular}{llrrrr}
\toprule
Bộ dữ liệu & Phương pháp & $T = 5$ & $T = 10$ & $T = 20$ & Không Phase 1 \\
\midrule
CIFAR100   & RanPAC & 92.4 & 92.2 & 90.8 & 89.0 \\
           & ADaM   & 88.5 & 87.5 & 85.2 & --- \\
\addlinespace
ImageNet-R & RanPAC & 79.9 & 77.9 & 74.5 & 71.8 \\
           & ADaM   & 74.3 & 72.9 & 70.5 & --- \\
\bottomrule
\end{tabular}
\end{table}

Đây là một trong những phát hiện có ích nhất của bài, và nó liên quan trực tiếp
tới đề tài của chúng ta.

Chia càng nhiều tác vụ thì Phase 1 càng mất tác dụng. Ở $T = 20$, RanPAC gần như
tụt về mức không có Phase 1 (ImageNet-R: 74.5 so với 71.8). Lý do rất rõ: chia
càng nhỏ thì tác vụ đầu tiên càng ít lớp, mà Phase 1 chỉ được huấn luyện trên
đúng tác vụ đó. Chính bài kết luận rằng chiến lược phiên đầu ``có thể ít giá trị
nếu tác vụ đầu tiên không đủ đa dạng''.

\begin{luuy}
Ràng buộc này áp dụng nguyên vẹn cho chúng ta. Vật thay thế Phase 1 trong hệ
thống của ta là bộ LoRA tác vụ 0 của HRM-PET, cũng chỉ được huấn luyện trên một
tác vụ duy nhất -- 20 lớp trên ImageNet-R. Nên không nên kỳ vọng nó mạnh bằng
Phase 1 thật.
\end{luuy}

\subsection{Tính bền vững}

Phụ lục F kiểm chứng thêm ba điều: phương pháp chạy được với backbone ResNet và
CLIP chứ không riêng ViT; nó áp dụng được cho học liên tục \emph{không có ranh
giới tác vụ} (task-agnostic, dùng giao thức Gaussian-scheduled CIFAR100); và kết
quả ổn định qua các seed, đặc biệt gần như bất biến khi bỏ Phase 1.

% =====================================================================
\section{Cái giá phải trả}

\begin{itemize}[leftmargin=*]
\item \textbf{Bộ nhớ khi huấn luyện.} Ma trận Gram kích thước $M \times M$, tức
      $10\,000^2$ phần tử. Đây là khoản lớn nhất. Bài lập luận rằng vẫn nhỏ hơn
      SLCA, vốn lưu một ma trận hiệp phương sai cho \emph{mỗi} lớp.
\item \textbf{Số tham số.} Thêm $L \times M = 7.68$ triệu trọng số đóng băng (có
      thể lưu 1 bit mỗi phần tử nếu dùng ma trận nhị cực), cùng khoảng 2--2.5
      triệu tham số học được với $K = 200$. So sánh: L2P và DualPrompt dùng
      khoảng 0.3--0.5 triệu.
\item \textbf{Tốc độ suy luận.} Gần như không đổi so với backbone gốc, vì lớp
      chiếu và đầu phân lớp đều chỉ là phép nhân ma trận.
\item \textbf{Tốc độ huấn luyện.} Phase 2 chỉ chậm hơn một lượt chạy suy luận
      chút ít. Phần chậm nhất là nghịch đảo ma trận Gram khi quét $\lambda$,
      khoảng 1 phút mỗi tác vụ trên CPU với $M = 10\,000$.
\item \textbf{Phần cứng.} Toàn bộ thí nghiệm của bài chạy trên một GPU
      RTX 4090 duy nhất.
\end{itemize}

% =====================================================================
\section{Hạn chế và hướng mở rộng}

\subsection{Hạn chế mà chính tác giả thừa nhận}

Toàn bộ giá trị của phương pháp \textbf{phụ thuộc hoàn toàn vào chất lượng của bộ
trích đặc trưng có sẵn}. Nó không học biểu diễn, nó chỉ khai thác biểu diễn. Vì
vậy nó khó có thể mạnh tương đương trong các thiết lập phải huấn luyện mạng từ
đầu.

Ngoài ra nó dùng nhiều tham số hơn các phương pháp prompting, dù tác giả cho rằng
đánh đổi này xứng đáng với sự đơn giản khi cài đặt.

\subsection{Hướng mở rộng tác giả đề xuất}

\begin{itemize}[leftmargin=*]
\item Học liên tục \textbf{không có ranh giới tác vụ}, tức là dữ liệu trôi dần
      chứ không chia thành đợt rõ ràng.
\item Dùng mục tiêu \textbf{không phải one-hot}, ví dụ nhúng ngôn ngữ của CLIP.
      Thí nghiệm sơ bộ cho $77.5\,\%$ với $M = 5000$, so với $68.6\,\%$ của CLIP
      zero-shot và $71.4\,\%$ khi không dùng chiếu.
\item Mở rộng sang \textbf{bài toán hồi quy}, khi đó khái niệm class-prototype
      phải tổng quát hóa thành feature-prototype.
\item Kết hợp với các phương pháp prompting, và thử huấn luyện PETL vượt quá
      phiên đầu tiên.
\end{itemize}

% =====================================================================
\section{Tóm tắt và liên hệ với đề tài của chúng ta}

\begin{luuy}
\textbf{Một câu tóm tắt cả bài:} RanPAC là \emph{hồi quy ridge chính xác, trên
đặc trưng đóng băng đã được chiếu ngẫu nhiên lên chiều cao kèm phi tuyến}. Và vì
nghiệm ridge chỉ phụ thuộc vào hai tổng cộng dồn, nó bất biến với thứ tự dữ liệu
-- nên bài toán quên thảm khốc không bị \emph{giải quyết}, mà bị định nghĩa cho
biến mất.
\end{luuy}

Trong đề tài của chúng ta, thành phần được lấy từ bài này là \textbf{toàn bộ
Phase 2}: lớp chiếu ngẫu nhiên đóng băng kèm ReLU, ma trận Gram, class-prototype
và nghiệm ridge. Chúng ta \textbf{không} cài Phase 1, mà thay bằng bộ LoRA tác vụ
0 mà HRM-PET vốn đã huấn luyện sẵn -- nhờ vậy không phải huấn luyện thêm một tham
số nào.

Bốn điểm cần nhớ khi đối chiếu số liệu:

\begin{itemize}[leftmargin=*]
\item RanPAC \textbf{không có bước định tuyến}. Nó là phương pháp một đầu phân
      lớp duy nhất trải trên mọi lớp; trong cả bài không có khái niệm suy luận
      danh tính tác vụ. Việc gộp điểm số theo lớp của nó thành điểm theo tác vụ
      để làm bộ định tuyến là cách dùng không có trong bài gốc.
\item Chỉ số $A_T$ của họ \textbf{trùng định nghĩa} với cột Acc@1 của chúng ta,
      nên hai bên so sánh được về mặt chỉ số. Tuy nhiên seed và cách chia lớp
      khác nhau, và chúng ta chưa từng tự chạy RanPAC, nên chưa thể coi là so
      sánh trực tiếp.
\item Vật thay thế Phase 1 của chúng ta chịu đúng ràng buộc nêu ở Mục
      \ref{sec:sotacvu}: nó chỉ thấy một tác vụ, nên khả năng vá khoảng cách miền
      bị giới hạn.
\item Bảng 1 và Bảng 2 của bài báo cho \textbf{số hơi khác nhau} cho cùng một
      phương pháp (ImageNet-R là 77.9 và 78.1; ImageNet-A là 62.4 và 61.8), nhiều
      khả năng vì Bảng 2 chọn backbone tốt nhất cho từng bộ dữ liệu. Khi trích số
      vào báo cáo cần ghi rõ lấy từ bảng nào.
\end{itemize}

\begin{thebibliography}{9}
\bibitem{ranpac}
Mark D. McDonnell, Dong Gong, Amin Parvaneh, Ehsan Abbasnejad, Anton van den
Hengel.
\newblock RanPAC: Random Projections and Pre-trained Models for Continual
Learning.
\newblock \emph{Advances in Neural Information Processing Systems 36 (NeurIPS
2023)}. arXiv:2307.02251.
\end{thebibliography}

\end{document}
